\chapter{Introduction}

\section{Background}

Personal Identification Numbers (PINs) are widely used for authentication in various systems, 
including banking, mobile devices, and online services. A 6-digit PIN provides a theoretical 
space of one million possible combinations. However, in practice, users tend to select PINs 
based on memorable patterns, such as dates, repeated digits, and sequences.

\section{Problem Statement}

The practical security of a PIN depends not only on its theoretical key space but also on how 
users choose their PINs. Human bias reduces the effective entropy of the system, making it 
vulnerable to guessing attacks. Furthermore, personal information such as birth dates can 
further increase the success rate of targeted attacks.

\section{Objectives}

The objective of this thesis is to evaluate the vulnerability of human-chosen 6-digit PINs 
under low-entropy attacks. The study aims to model user behavior, simulate attack strategies, 
and measure security using guessability-based metrics.

\section{Contributions}

This thesis provides:
\begin{itemize}
    \item A probabilistic model for human-chosen PIN distributions
    \item Simulation of multiple attack strategies
    \item Evaluation of security under limited guessing attempts
    \item Analysis of the impact of personal information leakage
\end{itemize}

\section{Structure of the Thesis}

The remainder of this thesis is organized as follows. Chapter 2 reviews related work on PIN 
security and low-entropy attacks. Chapter 3 presents the methodology. Chapter 4 discusses 
experimental results. Finally, Chapter 5 concludes the thesis.